% ============================================================================
% graficas/barras.tex
% Gráfico de barras por participante con línea de referencia global SUS (68)
% ============================================================================

\edef\susEtiquetas{\directlua{sus.imprimir_etiquetas_barras()}}
\edef\susCoordenadas{\directlua{sus.imprimir_coordenadas_barras()}}
\edef\susLineaRef{(\susPrimerId, 68) (\susUltimoId, 68)}

\begin{tikzpicture}
  \begin{axis}[
    ybar,
    width=0.95\linewidth,
    height=6.2cm,
    ymin=0,
    ymax=105,
    ylabel={Puntuación SUS},
    symbolic x coords/.expanded={\susEtiquetas},
    xtick=data,
    nodes near coords,
    nodes near coords style={font=\footnotesize\bfseries, text=susTexto},
    bar width=18pt,
    enlarge x limits=0.15,
    ymajorgrids=true,
    grid style={dashed, gray!30},
    axis line style={gray!60},
    tick label style={font=\small},
    ylabel style={font=\small\bfseries}
  ]
    % Barras con las puntuaciones de cada participante
    \addplot[
      fill=susPrincipal!85,
      draw=susPrincipal,
      line width=0.8pt
    ] coordinates {\susCoordenadas};

    % Línea de referencia estándar de la industria (SUS = 68)
    \addplot[
      color=susReferencia,
      dashed,
      line width=1.3pt
    ] coordinates {\susLineaRef};
    \node[susReferencia, fill=white, fill opacity=0.85, text opacity=1, inner sep=2pt, font=\scriptsize\bfseries] 
      at (rel axis cs:0.22, 0.69) {Media estándar de referencia (68)};
  \end{axis}
\end{tikzpicture}
